\documentclass[oneside,a4paper,14pt]{extarticle}
\usepackage[T1,T2A]{fontenc}
\usepackage[a4paper,letterpaper,top=20mm,bottom=20mm,left=30mm,right=10mm]{geometry}
\usepackage{amsmath, amssymb}
\usepackage[utf8]{inputenc}
\usepackage[russian]{babel}
\usepackage{textcomp}
\usepackage{indentfirst}
\usepackage{graphicx}
\usepackage{float}
\usepackage{mwe}
\usepackage{wrapfig}
\usepackage{caption}
\usepackage{tocloft}
\renewcommand{\cftsecleader}{\cftdotfill{\cftdotsep}}
\usepackage{longtable}
\usepackage{amsmath}
\usepackage{amsfonts}
\usepackage{amsthm}
\usepackage{float}
\usepackage[all]{xy}
\usepackage[breaklinks]{hyperref}
\usepackage{titlesec}
\usepackage{fancyvrb}
\usepackage{paracol}
\usepackage{listings}
\usepackage{xcolor}

\definecolor{codegreen}{rgb}{0,0.6,0}
\definecolor{codegray}{rgb}{0.5,0.5,0.5}
\definecolor{codepurple}{rgb}{0.58,0,0.82}
\definecolor{backcolour}{rgb}{0.95,0.95,0.92}

\lstdefinestyle{javastyle}{
    backgroundcolor=\color{backcolour},
    commentstyle=\color{codegreen},
    keywordstyle=\color{magenta},
    numberstyle=\tiny\color{codegray},
    stringstyle=\color{codepurple},
    basicstyle=\ttfamily\footnotesize,
    breakatwhitespace=false,
    breaklines=true,
    captionpos=b,
    keepspaces=true,
    numbers=left,
    numbersep=5pt,
    showspaces=false,
    showstringspaces=false,
    showtabs=false,
    tabsize=2,
    language=Java
}

\lstset{style=javastyle}

\titleformat{\section}{\normalsize\bfseries}{\thesection}{1em}{}
\titleformat{\subsection}{\normalsize\bfseries}{\thesubsection}{1em}{}
\titleformat{\subsubsection}{\normalsize\bfseries}{\thesubsubsection}{1em}{}
\renewcommand\baselinestretch{1.33}\normalsize
\setlength{\parindent}{1.25cm}

\begin{document}

\newpage\thispagestyle{empty}
\begin{center}
	МИНИСТЕРСТВО НАУКИ И ВЫСШЕГО ОБРАЗОВАНИЯ\\
	РОССИЙСКОЙ ФЕДЕРАЦИИ\\
	ФЕДЕРАЛЬНОЕ ГОСУДАРСТВЕННОЕ БЮДЖЕТНОЕ\\
	ОБРАЗОВАТЕЛЬНОЕ УЧРЕЖДЕНИЕ ВЫСШЕГО ОБРАЗОВАНИЯ\\
	«ВЯТСКИЙ ГОСУДАРСТВЕННЫЙ УНИВЕРСИТЕТ»\\
	Институт математики и информационных систем\\
	Факультет автоматики и вычислительной техники\\
	Кафедра электронных вычислительных машин
\end{center}

\vspace{20mm}
\begin{center}

	Отчёт по лабораторной работе №3\\
	по дисциплине\\
	<<Объектно-ориентированное программирование>>\\
\end{center}
\vfill

Выполнил студент гр. ИВТб-2301-05-00 \hfill
\makebox[8cm]{\hrulefill /Ковальский И.И./}

Проверил преподаватель\hfill
\makebox[8cm]{\hrulefill /Шмакова Н.А./}

\vfill
\begin{center}
	Киров\\
	2026
\end{center}

\newpage
\setcounter{page}{1}
\tableofcontents
\newpage

\section*{Цель работы}
\addcontentsline{toc}{section}{Цель работы}
Разработка приложения на языке Java с использованием объектно-ориентированного подхода, фреймворка Spring Boot и объектно-реляционного отображения (Hibernate) для хранения данных в реляционной базе данных PostgreSQL. Приложение должно реализовывать CRUD-операции над сущностями предметной области, поддерживать наследование и связи между таблицами в соответствии с принципами ООП и нормализацией баз данных.

\newpage

\section*{Задание}
\addcontentsline{toc}{section}{Задание}

На основании предметной области, согласованной в первой и второй лабораторных работах (тема «Организации»), разработать приложение на Java, удовлетворяющее следующим требованиям:

\begin{enumerate}
	\item Использование языка программирования Java.
	\item Использование СУБД PostgreSQL.
	\item Использование сборщика Maven (все зависимости подключаются через pom.xml).
	\item Использование Hibernate (JPA) для работы с БД — никаких SQL-запросов, только операции над сущностями.
	\item База данных должна содержать не менее трёх связанных таблиц, находящихся в третьей нормальной форме (3НФ).
	\item Для разметки сущностей использовать аннотации JPA, допускается применение Lombok для сокращения кода.
	\item Для реализации CRUD-операций использовать интерфейсы, наследующие \texttt{CrudRepository} из Spring Data.
	\item Сохранить структуру классов из предыдущих работ: абстрактный базовый класс, наследование, инкапсуляция полей через геттеры/сеттеры.
	\item Обеспечить возможность вызова методов классов (отчёт, налоги, уникальные методы) через пользовательский интерфейс.
	\item В отчёте описать диаграмму классов, включая классы-сущности, репозитории и контроллеры, а также особенности подключения к БД.
\end{enumerate}

\newpage

\subsection*{Описание предметной области и иерархии классов}
\addcontentsline{toc}{subsection}{Описание предметной области и иерархии классов}

Предметная область — «Организации». В третьей лабораторной работе модель была расширена: добавлена сущность \texttt{Employees} (сотрудники), которая связана с организацией отношением «многие к одному». Таким образом, база данных содержит три основные таблицы: \texttt{organizations} (базовая), \texttt{comorganizations} и \texttt{noncomorganizations} (дочерние, для разных типов организаций), а также таблицу \texttt{employees}. Все таблицы находятся в третьей нормальной форме.

\paragraph{Иерархия классов-сущностей (JPA):}
\begin{itemize}
	\item \texttt{Organization} — абстрактный базовый класс, аннотированный как \texttt{@Entity} со стратегией наследования \texttt{InheritanceType.JOINED}. Содержит общие поля: \texttt{id}, \texttt{name}, \texttt{inn}, \texttt{employeesCount} и связь \texttt{@OneToMany} со списком сотрудников.
	\item \texttt{ComOrg} — дочерний класс коммерческой организации. Добавляет поля \texttt{profit} (прибыль) и \texttt{businessType} (тип бизнеса). Реализует абстрактные методы \texttt{report()} и \texttt{taxes()}, а также собственные методы \texttt{expandBusiness()} и \texttt{distribute()}.
	\item \texttt{NonComOrg} — дочерний класс некоммерческой организации. Добавляет поля \texttt{purpose} (цель) и \texttt{source} (источник финансирования). Переопределяет методы базового класса, содержит уникальные методы \texttt{program()} и \texttt{attractFunding()}.
	\item \texttt{Employees} — самостоятельная сущность, представляющая сотрудника. Содержит поля \texttt{id}, \texttt{name}, \texttt{position} и ссылку \texttt{@ManyToOne} на организацию.
\end{itemize}

\newpage

\paragraph{Слои приложения:}
\begin{itemize}

	\item \textbf{Сущности} — описанные выше классы с аннотациями JPA.

    \item \textbf{Репозитории} — интерфейсы, расширяющие \texttt{CrudRepository} (например, \texttt{OrganizationRep}, \texttt{ComOrgRep}, \texttt{NonComOrgRep}, \texttt{EmployeesRep}). Они предоставляют готовые методы для CRUD-операций без явной реализации.

    \item \textbf{Контроллер} — класс \texttt{WebController}, аннотированный \texttt{@Controller}. Обрабатывает HTTP-запросы, вызывает методы репозиториев и возвращает HTML-страницы (представления Thymeleaf).

    \item \textbf{Точка входа} — класс \texttt{App} с аннотацией \texttt{@SpringBootApplication}, запускающий встроенный веб-сервер.

\end{itemize}

% \begin{figure}[H]
%   \includegraphics[width=1\textwidth]{Диаграмма классов.png}
%   \caption{Диаграмма классов}
%   \label{fig:screen form}
% \end{figure}

\newpage

\subsection*{Особенности подключения и работы с PostgreSQL через Hibernate}
\addcontentsline{toc}{subsection}{Особенности подключения и работы с БД PostgreSQL через Hibernate}

В третьей лабораторной работе для доступа к базе данных используется Hibernate в качестве реализации JPA, а также Spring Data JPA, который автоматически генерирует реализации репозиториев. Конфигурация подключения задаётся в файле \texttt{application.properties} (или \texttt{application.yml}):

\begin{lstlisting}[caption=application.properties, label=lst:props]
spring.datasource.url=jdbc:postgresql://localhost:5432/lab3_db
spring.datasource.username=postgres
spring.datasource.password=****
spring.datasource.driver-class-name=org.postgresql.Driver

spring.jpa.hibernate.ddl-auto=update
spring.jpa.show-sql=true
spring.jpa.properties.hibernate.dialect=org.hibernate.dialect.PostgreSQLDialect
spring.jpa.properties.hibernate.format_sql=true
\end{lstlisting}

Параметр \texttt{ddl-auto=update} позволяет Hibernate автоматически создавать или обновлять схему базы данных на основе аннотаций сущностей, что удобно при разработке.

\paragraph{Особенности реализации:}
\begin{itemize}
	\item \textbf{Стратегия наследования} — \texttt{JOINED}. Для каждого класса в иерархии создаётся отдельная таблица, а дочерние таблицы связаны с родительской через первичный ключ (внешний ключ). Это обеспечивает нормализацию и соответствует требованию «не хранить данные скопом в одной таблице».
	\item \textbf{Связи между таблицами}: таблица \texttt{employees} имеет внешний ключ \texttt{org\_id}, ссылающийся на \texttt{organizations.id}. В сущностях это выражено аннотациями \texttt{@OneToMany} и \texttt{@ManyToOne} с каскадными операциями.
	\item \textbf{Использование Lombok}: аннотации \texttt{@Data}, \texttt{@NoArgsConstructor}, \texttt{@AllArgsConstructor}, \texttt{@EqualsAndHashCode} автоматически генерируют конструкторы, геттеры, сеттеры и методы \texttt{equals/hashCode}, сокращая объём шаблонного кода.
	\item \textbf{Репозитории}: наследуя \texttt{CrudRepository}, мы получаем готовые методы \texttt{save()}, \texttt{findById()}, \texttt{findAll()}, \texttt{deleteById()} и т.д. Это полностью исключает написание SQL-запросов.
	\item \textbf{Транзакционность}: методы контроллера неявно выполняются в транзакциях благодаря аннотации \texttt{@Transactional} на репозиториях (обеспечивается Spring Data JPA).
\end{itemize}

Таким образом, приложение полностью соответствует требованию «никаких SQL-запросов в коде» — вся работа с БД производится через манипуляции с объектами Java.

\newpage

\subsection*{Пользовательского web-интерфейс}
\addcontentsline{toc}{subsection}{Разработка пользовательского интерфейса}

В третьей работе был реализован веб-интерфейс с использованием Spring MVC. Пользователь взаимодействует с приложением через браузер.\\

Главная страница содержит:

\begin{itemize}
	\item Список всех организаций, полученный через \texttt{organizationRep.findAll()}.
	\item Список всех сотрудников.
	\item Формы для добавления новой коммерческой/некоммерческой организации.
	\item Формы для добавления нового сотрудника с возможностью привязки к существующей организации.
	\item Для каждой организации и каждого сотрудника доступны кнопки «Редактировать» и «Удалить», которые ведут на отдельные страницы или модальные окна | POST-запросы.
\end{itemize}


Контроллер \texttt{WebController} обрабатывает GET-запросы для отображения страниц и POST-запросы для выполнения операций создания, обновления и удаления. Например, метод \texttt{addComOrg()} принимает параметры формы, создаёт объект \texttt{ComOrg} и сохраняет его через \texttt{comOrgRep.save()}. После выполнения всех операций происходит перенаправление на главную страницу.

Для вызова методов предметной области (например, \texttt{expandBusiness()}, \texttt{taxes()}, \texttt{report()}) в интерфейсе предусмотрены отдельные кнопки или ссылки, которые отправляют запрос на контроллер, вызывающий соответствующий метод сущности и возвращающий результат в виде отдельной страницы или всплывающего сообщения (в реализованном коде такие методы вызываются непосредственно в контроллере, а результат отображается через модель).

Преимущества такого подхода:
\begin{itemize}
	\item Кроссплатформенность — интерфейс доступен с любого устройства через браузер.
	\item Чёткое разделение логики: модель (сущности), представление (HTML-шаблоны), контроллер (обработка запросов).
	\item Лёгкость масштабирования и добавления новых функций.
\end{itemize}

\newpage

\section*{Вывод}
\addcontentsline{toc}{section}{Вывод}

В ходе выполнения третьей лабораторной работы было разработано веб-приложение на Java с использованием Spring Boot, Spring Data JPA (Hibernate) и PostgreSQL. Приложение полностью удовлетворяет поставленным требованиям:

\begin{enumerate}
	\item Сохранена объектно-ориентированная структура классов: абстрактный базовый класс \texttt{Organization}, дочерние классы \texttt{ComOrg} и \texttt{NonComOrg}, инкапсуляция полей через геттеры/сеттеры (сгенерированные Lombok).
	\item Для хранения данных использована СУБД PostgreSQL, спроектированы три связанные таблицы (\texttt{organizations}, \texttt{comorganizations}, \texttt{noncomorganizations} и дополнительная \texttt{employees}), что обеспечивает третью нормальную форму.
	\item Реализованы все CRUD-операции через репозитории Spring Data, без написания явных SQL-запросов.
	\item Взаимодействие с БД построено на Hibernate (JPA) с использованием аннотаций и стратегии наследования \texttt{JOINED}.
	\item Пользовательский интерфейс выполнен в виде веб-приложения с шаблонами Thymeleaf, что позволяет удобно управлять данными и вызывать методы классов.
\end{enumerate}

Таким образом, полученный проект демонстрирует навыки построения полноценного enterprise-приложения с современным стеком технологий (Java, Spring, Hibernate, PostgreSQL, Maven), а также умение применять принципы ООП и нормализации баз данных.

\end{document}